% !TEX program = xelatex
% ============================================================
%  CUMCM 论文 LaTeX 骨架（主文件模板）
%  编译：latexmk -xelatex main.tex   （或双击 build.bat）
%  规范要点：电子版首页 = 摘要页；不含承诺书/编号页；正文 <=30 页；不加目录
%  用法：复制本目录（论文骨架）为你的论文工程 → 填 sections/ 各章
%        → 摘要页最后写；编译用 build.bat 或 compile_check.sh
% ============================================================
\documentclass[12pt, a4paper]{ctexart}
% ============================================================
%  导言区：样式与宏包（CUMCM 论文通用模板，实测可编译）
%  依据：官方《全国大学生数学建模竞赛论文格式规范》
%        —— 页边距上下左右 >=2.5cm；字体/字号/行距官方未强制
% ============================================================

% ---------- 页面 ----------
\usepackage{geometry}
\geometry{a4paper, top=2.5cm, bottom=2.5cm, left=2.5cm, right=2.5cm}

% ---------- 行距（为满足正文 <=30 页限制，由 1.5 调整为 1.38） ----------
\linespread{1.38}

% ---------- 页码：页脚居中（官方规范） ----------
\pagestyle{plain}

% ---------- 章节编号：中文数字 + 、（如 “一、问题重述”） ----------
\ctexset{
  section = {number = \chinese{section}, aftername = {、}},
}

% ---------- 字体：正文中文宋体（ctex 默认）；西文/数字/公式为 Times（MathType 观感） ----------
\setmainfont{Times New Roman}

% ---------- 数学 ----------
\usepackage{amsmath, amsfonts, mathtools}
\usepackage{unicode-math}
\setmathfont{XITSMath-Regular.otf}   % XITS Math = Times 系 OpenType 数学字体

% ---------- 定理环境 ----------
\usepackage{amsthm}
\newtheorem{lemma}{引理}
\newtheorem{theorem}{定理}
\newtheorem{proposition}{命题}
\newtheorem{corollary}{推论}

% ---------- 图表 ----------
\usepackage{graphicx}
\graphicspath{{figures/}}
\usepackage{booktabs}
\usepackage{multirow}
\usepackage{makecell}
\usepackage{threeparttable}
\usepackage{tabularx}
\usepackage{longtable}
\usepackage{float}
\usepackage{caption}
\usepackage{subcaption}
\captionsetup{font=small, labelsep=quad}

% ---------- 列表 ----------
\usepackage{enumitem}

% ---------- 颜色与代码 ----------
\usepackage{xcolor}
\usepackage{listings}
\lstset{
  basicstyle=\ttfamily\small,
  breaklines=true,
  frame=single,
  numbers=left,
  numberstyle=\tiny,
}

% ---------- 算法伪代码 ----------
\usepackage[ruled, vlined, linesnumbered]{algorithm2e}
\SetAlgorithmName{算法}{算法}{算法索引}

% ---------- 单位 ----------
\usepackage{siunitx}

% ---------- 参考文献（GB/T 7714—2015 顺序编码制） ----------
\usepackage{natbib}
\usepackage{gbt7714}

% ---------- 超链接（放在最后加载） ----------
\usepackage[colorlinks=true, linkcolor=black, citecolor=black, urlcolor=blue]{hyperref}
\usepackage{xurl}

% ---------- 浮动体版面参数（减少整页空白） ----------
\renewcommand{\topfraction}{0.9}
\renewcommand{\bottomfraction}{0.85}
\renewcommand{\textfraction}{0.07}
\renewcommand{\floatpagefraction}{0.72}
\setcounter{topnumber}{3}
\setcounter{bottomnumber}{2}
\setcounter{totalnumber}{4}


\begin{document}

% ---------- 摘要页（页码从 1 开始） ----------
% ============================================================
%  摘要页（电子版第一页；页码从此页开始，为“1”）
%  红线：一页装下；最后写；数字写实；关键词用“文献高频常用词”
% ============================================================
\thispagestyle{plain}

\begin{center}
  {\zihao{3}\heiti 论文标题（三号黑体居中）\par}
  \vspace{0.1em}
  {\zihao{4}\heiti 摘\quad 要\par}
\end{center}
\vspace{0.08em}
{\linespread{1.20}\selectfont

总起段：面向什么问题、建立什么样的“递进模型”，并用什么链条验证（1 段）。

针对问题一，……（方法 → 关键数值结果 3~6 位 → 验证口径一句话）。
针对问题二，……
针对问题三，……
针对问题四，……

特色总结段（1 段）：方法亮点 + 验证链完整性 + 对局限/近似的如实说明。

\vspace{0.05em}
\noindent{\heiti 关键词：}\quad 词1；词2；词3；词4；词5\par}
\newpage


% ---------- 正文（官方要求不加目录） ----------
% 一、问题重述：用自己的话重述背景与任务（不抄题；每问一句话点题）
\section{问题重述}

（1）背景：……

（2）任务：问题一要求……；问题二要求……；问题三要求……；问题四要求……。

% 二、问题分析：每问写清“已知/未知/难点/思路”，并标出“可验证点”（评委看什么）
\section{问题分析}

2.1 问题一分析：……

2.2 问题二分析：……

2.3 问题三分析：……

2.4 问题四分析：……

% 三、模型假设：逐条编号，写明依据；宁少勿滥（多余假设会削弱模型）
\section{模型假设}

假设 1：……（依据：……）

假设 2：……（依据：……）

假设 3：……（依据：……）

% 四、符号说明：表格（符号 / 含义 / 单位）；全文符号必须统一，表格只列核心符号
\section{符号说明}

\begin{table}[htbp]
  \centering
  \small
  \begin{tabular}{clc}
    \toprule
    符号 & 含义 & 单位 \\
    \midrule
    $x$ & …… & m \\
    \bottomrule
  \end{tabular}
\end{table}

% 问题一：模型建立 → 求解 → 结果（给数字）→ 分析与验证小节（评委必看）
% 提醒：所有解析结论先过数值验证（交叉验证 + 批量统计 + 灵敏度），见手册 §5
\section{问题一：模型的建立与求解}

5.1 模型建立：……

5.2 求解方法：……

5.3 结果与分析：……

5.4 验证：……（交叉验证 / 批量统计 / 灵敏度，给出具体数字）

% 问题二：模型建立 → 求解 → 结果（给数字）→ 分析与验证小节
% 多方案时用“选择—对比—取舍”写法：列备选 → 统一判据对比 → 说明取舍理由
\section{问题二：模型的建立与求解}

6.1 模型建立：……

6.2 求解方法：……

6.3 结果与分析：……

6.4 验证：……

% 问题三：策略设计 → 参数优化 → 批量演练统计 → 结果数字 → 验证与审计
% 要点：多方案对比取舍（给统一判据）；批量演练给清除率/时间分布；动作计时逐条审计
\section{问题三：模型的建立与求解}

7.1 策略设计：……

7.2 参数优化：……

7.3 批量演练与结果：……（清除比例、平均时间等具体数字）

7.4 验证与审计：……

% 问题四：扩展/不确定性类（问题辨析 → 判据升级 → 定量刻画 → 策略修正 → 验证）
% 要点：把新情况带来的“歧义”逐条显式列出并给判据（如“几重锁”）；风险定量化；后备方案与可验证证书
\section{问题四：模型的建立与求解}

8.1 问题辨析与判据设计：……

8.2 风险定量刻画：……

8.3 策略修正：……

8.4 验证：……

% 九、模型评价与推广：优点 / 缺点与改进方向 / 推广应用
% 要点：缺点少而缓（2~3 条，说清改进路径）；优点对应评委关注点；推广给具体场景
\section{模型评价与推广}

9.1 模型优点：……

9.2 模型缺点与改进方向：……

9.3 推广应用：……


% ---------- AI 工具使用声明（按当届规定：置于参考文献之前） ----------
\section*{AI 工具使用声明}
本参赛队在竞赛过程中使用了 AI 工具，主要用于辅助程序开发与调试、数据整理与思路探讨，详细使用情况见下。

\begin{itemize}[leftmargin=2em, itemsep=2pt, topsep=3pt]
  \item 工具A（名称/版本）：用途环节（示例：算法程序纠错与优化）；
  \item 工具B（名称/版本）：用途环节（示例：文档整理与配图）。
\end{itemize}

% ---------- 参考文献（GB/T 7714—2015，顺序编码制） ----------
\clearpage
\bibliographystyle{gbt7714-2015-numeric}
\bibliography{refs}

% ---------- 附录（支撑材料目录 + 运行说明 + 核心程序清单） ----------
\clearpage
% 附录：A 支撑材料目录（条目让评委一眼看到“主程序”）/ B 运行说明 / C 核心程序清单（节选）
% 完整源码放支撑材料承载；附录 C 只截取核心问题的核心代码（思想核心、少而精）
% 附录条目命名建议直接体现身份：如 01_主程序_xxx、主程序_xxx.py

\section*{附录 A：支撑材料目录}

（示例条目：01_主程序_xxx / 02_问题求解程序 / 03_验证与工具 / 04_统计数据 / 05_正式日志）

\section*{附录 B：程序运行说明}

（环境要求；运行命令示例：\texttt{python main.py}；依赖说明。）

\section*{附录 C：核心程序清单（节选）}

（只放核心函数 / 主流程节选，每段前一句话说明作用。）


\end{document}
